\documentclass[12pt,a4paper,oneside]{report}
\usepackage[utf8]{inputenc}
\usepackage[T1]{fontenc}
\usepackage[portuguese]{babel}
\usepackage[left=2.5cm,right=2.5cm,top=2.5cm,bottom=2.5cm]{geometry}
\usepackage{graphicx}
\usepackage[table]{xcolor}
\usepackage{array}
\usepackage{booktabs}
\usepackage{longtable}
\usepackage{float}
\usepackage{fancyhdr}
\usepackage{hyperref}
\usepackage{setspace}
\usepackage{tabularx}
\usepackage{tikz}
\usetikzlibrary{arrows.meta,positioning,shapes.geometric,calc}
\usepackage{enumitem}
\usepackage{lmodern}
\usepackage{microtype}
\usepackage{listings}

\definecolor{primary}{RGB}{18,56,102}
\definecolor{secondary}{RGB}{28,120,86}
\definecolor{lightblue}{RGB}{236,242,250}
\definecolor{lightgreen}{RGB}{235,247,241}
\definecolor{graytext}{RGB}{80,80,80}
\definecolor{linecolor}{RGB}{80,110,150}
\definecolor{codebackground}{RGB}{245,245,245}

\hypersetup{
    colorlinks=true,
    linkcolor=primary,
    urlcolor=secondary,
    pdftitle={Projeto de Sistemas Distribuídos - Etapa 1},
    pdfauthor={Leonardo Cardoso de Moura}
}

\onehalfspacing
\pagestyle{fancy}
\fancyhf{}
\rhead{\thepage}
\lhead{Etapa 1 - Concepção e Arquitetura Inicial}
\renewcommand{\headrulewidth}{0.4pt}

\setlength{\parindent}{0pt}
\setlength{\parskip}{0.6em}

\lstset{
    language=C,
    basicstyle=\ttfamily\small,
    backgroundcolor=\color{codebackground},
    breaklines=true,
    frame=single,
    rulecolor=\color{primary},
    showstringspaces=false,
    keywordstyle=\color{primary},
    commentstyle=\color{gray},
    stringstyle=\color{secondary}
}

\title{
    \vspace{-1.5cm}
    {\Large \textbf{TRABALHO FINAL - ETAPA 1}}\\
    \vspace{0.4cm}
    {\LARGE \textbf{Concepção e Arquitetura Inicial}}\\
    \vspace{0.8cm}
    {\Large \textbf{Painel IoT Distribuído com Monitoramento em Tempo Real}}\\
    \vspace{0.8cm}
    {\large Infraestrutura Heterogênea: Raspberry Pi, Linux Mint e ESP32/ESP8266}\\
    \vspace{1.2cm}
}

\author{
    \textbf{Leonardo Cardoso de Moura}\\
    \textit{Matrícula: 32021BSI004}\\
    \vspace{0.3cm}
    Disciplina: Sistemas Distribuídos\\
    Data de entrega: 08 de maio de 2026\\
    Apresentação: 11 de maio de 2026
}

\date{}

\begin{document}

\maketitle
\thispagestyle{empty}
\newpage
\tableofcontents
\newpage

\chapter{Informações Gerais}

\begin{center}
\rowcolors{2}{lightblue}{white}
\begin{tabular}{|p{5cm}|p{9cm}|}
\hline
\rowcolor{primary}
\color{white}\textbf{Campo} & \color{white}\textbf{Valor} \\
\hline
Nome completo & Leonardo Cardoso de Moura \\
\hline
Matrícula & 32021BSI004 \\
\hline
Curso & Sistemas de Informação \\
\hline
Disciplina & Sistemas Distribuídos \\
\hline
Título do projeto & Painel IoT Distribuído com Monitoramento em Tempo Real \\
\hline
\end{tabular}
\end{center}

\chapter{Tema do Micro-Mundo}

\section{Descrição geral}
O micro-mundo proposto é um \textbf{painel IoT distribuído heterogêneo} para monitoramento em tempo real de ambiente, implementado com uma infraestrutura real composta por: Raspberry Pi 3 Model B em Linux Debian CLI funcionando como servidor de banco de dados, um notebook com Linux Mint hospedando a API e o sistema completo com gateway HTTP, múltiplos nós ESP32 e ESP8266 para aquisição de dados via sensores, e comunicação intermediária através de Nginx.

Os nós sensores (ESP32/ESP8266) não estarão concentrados em uma única residência; serão distribuídos em diferentes locais da cidade ou em quaisquer locais com conexão à Internet. O sistema simula um cenário genuinamente distribuído onde dispositivos geograficamente dispersos coletam dados, enfrentam qualidade de conexão variável, exibem informações localmente e sincronizam estado com o servidor central.

\section{O que o sistema faz}
\begin{enumerate}[leftmargin=1.2cm]
    \item Coleta contínua de leituras de sensores (temperatura, umidade, pressão).
    \item Monitoramento de força e qualidade de sinal WiFi em cada nó para diagnóstico de cobertura.
    \item Exibição em tempo real em telas OLED com animações responsivas.
    \item Controle interativo via sensor de toque capacitivo (mudança de expressão).
    \item Sincronização com servidor central através de API HTTP via Nginx.
    \item Persistência de séries temporais em banco de dados PostgreSQL no Raspberry Pi.
    \item Dashboard web para visualização, alertas e controle remoto.
\end{enumerate}

\section{Entidades do micro-mundo}
As entidades principais são:
\begin{itemize}[leftmargin=1.2cm]
    \item \textbf{Raspberry Pi 3 Model B}: servidor de banco de dados persistente.
    \item \textbf{Notebook Linux Mint}: servidor API e gateway HTTP/Nginx.
    \item \textbf{Múltiplos ESP32/ESP8266}: nós sensores distribuídos geograficamente (cidades, locais com conexão à Internet).
    \item \textbf{Sensores}: BMP280 (temperatura, pressão) em alguns nós e DHT11 (temperatura, umidade) em outros; não há redundância de sensores planejada.
    \item \textbf{Displays OLED}: exibição local de informações.
    \item \textbf{Sensor Toque TTP223}: entrada interativa para animação de rosto.
    \item \textbf{Nginx}: proxy reverso e balanceador de carga na rede local.
    \item \textbf{PostgreSQL}: banco de dados de séries temporais no Raspberry Pi.
\end{itemize}

\section{Interação central}
A interação principal do sistema é:
\begin{quote}
ESP32/ESP8266 $\rightarrow$ [WiFi] $\rightarrow$ Nginx (Notebook) $\rightarrow$ API (Node.js) $\rightarrow$ PostgreSQL (Raspberry Pi)
\end{quote}

Esse fluxo é interessante porque introduz desafios reais de sistemas distribuídos: latência de rede doméstica, falhas de conectividade WiFi, sincronização entre nós com qualidade de sinal variável, e consistência de estado compartilhado em um banco de dados remoto.

\chapter{Infraestrutura Detalhada}

\section{Servidores}

\subsection{Raspberry Pi 3 Model B - Servidor de Banco de Dados}
\begin{itemize}[leftmargin=1.2cm]
    \item \textbf{Sistema Operacional}: Linux Debian CLI (sem interface gráfica)
    \item \textbf{Função Principal}: PostgreSQL, armazenamento persistente, histórico de dados
    \item \textbf{Conectividade}: Ethernet cabeada para estabilidade máxima
    \item \textbf{IP Estático}: Fixado no roteador via DHCP com MAC binding (ex: 192.168.1.10)
    \item \textbf{Porta de Banco}: 5432 (PostgreSQL, acessível apenas do Notebook)
    \item \textbf{Benéfício Distribuído}: separação de responsabilidade, persistência centralizada, escalabilidade horizontal futura
    \item \textbf{Uptime}: 99,5\% (manutenção preventiva semanal)
\end{itemize}

\subsection{Notebook - Servidor de API e Gateway}
\begin{itemize}[leftmargin=1.2cm]
    \item \textbf{Sistema Operacional}: Linux Mint
    \item \textbf{Função Principal}: Node.js API, processamento de telemetria, dashboard web, gerenciamento de sessão
    \item \textbf{Conectividade}: WiFi doméstica
    \item \textbf{IP Estático}: Fixado no roteador via DHCP com MAC binding (ex: 192.168.1.11)
    \item \textbf{Portas de Serviço}:
    \begin{itemize}
        \item 3000: Node.js API backend
        \item 3001: Gateway de ingestão de telemetria
        \item 8080: Nginx proxy reverso
    \end{itemize}
    \item \textbf{Benéfício Distribuído}: separação entre ingestão de dados e armazenamento, ponto centralizado de lógica de negócio
    \item \textbf{Papel Futuro}: load balancing para múltiplos servidores API em fases posteriores
\end{itemize}

\section{Configuração de Rede Doméstica}

\subsection{DHCP Estático e IP Fixo}
Para garantir acesso remoto confiável e comunicação consistente entre nós, foi implementada estratégia de IP estático no roteador:

\begin{itemize}[leftmargin=1.2cm]
    \item \textbf{Problema}: internet residencial com IP dinâmico do provedor muda periodicamente, causando interrupção no acesso remoto
    \item \textbf{Solução}: fixação de endereços MAC dos servidores no roteador, atribuição de IPs privados permanentes via DHCP binding
    \item \textbf{Raspberry Pi}: \texttt{192.168.1.10} (fixo via MAC binding)
    \item \textbf{Notebook}: \texttt{192.168.1.11} (fixo via MAC binding)
    \item \textbf{ESP32/ESP8266}: IPs dinâmicos com prefixo reservado no roteador (DHCP range contíguo entre .100-.110)
    \item \textbf{Benefício}: eliminação de necessidade de DNS dinâmico ou discovery automático dentro da rede local
\end{itemize}

\subsection{Nginx como Proxy Reverso e Gateway}
\begin{itemize}[leftmargin=1.2cm]
    \item \textbf{Localização}: Notebook (Linux Mint)
    \item \textbf{Portas Públicas}: 8080 (HTTP local), 80 (futuro HTTPS)
    \item \textbf{Responsabilidades}:
    \begin{itemize}
        \item Roteamento de requisições HTTP POST dos ESP's para API (porta 3000)
        \item Compressão gzip de respostas para reduzir uso de largura de banda
        \item Balanceamento de carga para requisições futuras (múltiplas instâncias de API)
        \item Tratamento de CORS para acesso de clientes JavaScript nos ESP's
        \item Rate limiting para proteção contra abuso ou erros de firmware
        \item Logging de todas as requisições para auditoria
    \end{itemize}
    \item \textbf{Arquitetura}: single point of entry para todos os ESP's, evita comunicação direta com API
\end{itemize}

\subsection{Acesso Remoto Seguro}
Embora ainda em planejamento para futuras etapas:
\begin{itemize}[leftmargin=1.2cm]
    \item \textbf{Tailscale}: VPN pessoal para acesso de fora da rede local sem exposição de portas
    \item \textbf{Funnel (futuro)}: exposição temporária de endpoints para demonstrações
    \item \textbf{SSH}: acesso administrativo remoto do Raspberry Pi com chave pública
    \item \textbf{Firewall}: apenas portas 8080 e SSH expostas (futura regulação)
\end{itemize}

\chapter{Nós Sensores - ESP32 e ESP8266}

\section{Características Gerais}

Os nós sensores distribuídos são baseados em microcontroladores WiFi de baixo custo:

\begin{itemize}[leftmargin=1.2cm]
    \item \textbf{Quantidade}: múltiplos nós (3-5) espalhados pela cidade ou em diferentes locais com conexão à Internet
    \item \textbf{Modelos}: 
    \begin{itemize}
        \item ESP32: mais potente, dual-core, 4MB SPIFFS, maior consumo
        \item ESP8266: mais simples e econômico, single-core, 4MB SPIFFS
    \end{itemize}
    \item \textbf{Comunicação}: WiFi 802.11b/g/n para servidor API via Nginx
    \item \textbf{Protocolo de Dados}: HTTP POST com payloads JSON
    \item \textbf{Frequência de Amostragem}: 30-60 segundos (configurável via EEPROM)
    \item \textbf{Buffer Local}: retenção de últimas 10 leituras em caso de desconexão WiFi
    \item \textbf{Auto-recover}: reconexão automática com backoff exponencial
\end{itemize}

\section{Sensores Instalados}

\subsection{BMP280 - Sensor de Temperatura e Pressão}
\begin{itemize}[leftmargin=1.2cm]
    \item \textbf{Medidas Fornecidas}: 
    \begin{itemize}
        \item Temperatura: -40 a +85°C
        \item Pressão atmosférica: 300 a 1100 hPa
    \end{itemize}
    \item \textbf{Protocolo de Comunicação}: I2C (endereço 0x76 ou 0x77) ou SPI
    \item \textbf{Precisão}: $\pm 1.0°C$ (temperatura), $\pm 1 \text{ hPa}$ (pressão)
    \item \textbf{Tempo de Resposta}: < 1 segundo
    \item \textbf{Uso}: aquisição de temperatura e pressão ambiental em diversos nós
    \item \textbf{Nós}: alguns ESP utilizarão BMP280; outros utilizarão DHT11 conforme custo/disponibilidade
    \item \textbf{Vantagem Distribuída}: boa precisão para temperatura e pressão, baixo custo e fácil integração
\end{itemize}

\subsection{DHT11 - Sensor Temperatura/Umidade}
\begin{itemize}[leftmargin=1.2cm]
    \item \textbf{Medidas Fornecidas}:
    \begin{itemize}
        \item Temperatura: 0 a 50°C
        \item Umidade relativa: 20 a 80\%
    \end{itemize}
    \item \textbf{Protocolo de Comunicação}: DHT (single-wire digital proprietário)
    \item \textbf{Precisão}: $\pm 2°C$ (temperatura), $\pm 5\%$ (umidade)
    \item \textbf{Frequência de Amostragem}: máximo uma leitura a cada 1 segundo
    \item \textbf{Uso}: medição de temperatura e umidade em alguns nós; alguns dispositivos usarão DHT11 por custo ou disponibilidade
    \item \textbf{Nós}: alguns ESP (não todos, para economia de pinos e energia)
    \item \textbf{Benefício Distribuído}: baixo custo e simplicidade em locais onde apenas temperatura/umidade são necessárias
\end{itemize}

\section{Monitoramento de Qualidade de Sinal WiFi - RSSI}

\subsection{Received Signal Strength Indicator (RSSI)}
\begin{itemize}[leftmargin=1.2cm]
    \item \textbf{O que mede}: força do sinal WiFi em dBm (valores negativos)
    \item \textbf{Faixa típica}:
    \begin{itemize}
        \item -30 dBm ou melhor: excelente, máxima taxa de transferência
        \item -45 dBm: muito bom, confiável
        \item -60 dBm: aceitável, algumas quedas esperadas
        \item -80 dBm: fraco, desconexões frequentes
        \item -90 dBm ou pior: praticamente inutilizável
    \end{itemize}
    \item \textbf{Coleta}: a cada ciclo de amostragem, junto com temperatura/umidade
    \item \textbf{Propósito Distribuído}: 
    \begin{itemize}
        \item Diagnóstico real de cobertura WiFi em cada cômodo
        \item Detecção automática de interferência de frequência (outras redes, micro-ondas, Bluetooth)
        \item Planejamento futuro de repetidores ou mudança de canal
        \item Trigger de alertas se cobertura fica crítica
    \end{itemize}
    \item \textbf{Persistência}: armazenado no banco de dados para análise histórica e gráficos
    \item \textbf{Correlação}: possibilidade de correlacionar RSSI com taxa de erro de transmissão
\end{itemize}

\section{Telas OLED e Exibição Local}

\subsection{Displays OLED 128x64 - SSD1306}
\begin{itemize}[leftmargin=1.2cm]
    \item \textbf{Localização}: alguns ESP32 (não todos, aqueles em áreas visíveis)
    \item \textbf{Protocolo de Comunicação}: I2C (endereço 0x3C ou 0x3D)
    \item \textbf{Resolução}: 128 pixels (horizontal) × 64 pixels (vertical)
    \item \textbf{Cores}: monocromático (branco/preto)
    \item \textbf{Conteúdo Exibido}:
    \begin{itemize}
        \item Temperatura e umidade em tempo real (atualizados a cada amostragem)
        \item Força de sinal WiFi em barras ou valor em dBm
        \item Animação de rosto com expressões (animação em pixel art)
        \item Informação de status (``CONECTADO'', ``ENVIANDO'', ``ERRO'')
        \item Horário sincronizado via NTP
    \end{itemize}
    \item \textbf{Taxa de Atualização}: 1-2 Hz (suficiente para animação suave)
    \item \textbf{Consumo de Energia}: mínimo, não afeta bateria
\end{itemize}

\subsection{Sensor de Toque TTP223 - Interatividade Local}
\begin{itemize}[leftmargin=1.2cm]
    \item \textbf{Função}: controle por toque capacitivo sem mecânica física
    \item \textbf{Entrada}: GPIO digital (nível HIGH quando tocado, LOW quando solto)
    \item \textbf{Sensibilidade}: ajustável via solder jumper
    \item \textbf{Integração com OLED}: botão tátil para mudar expressão do rosto em tempo real
    \item \textbf{Expressões Implementadas}:
    \begin{itemize}
        \item Feliz: :) com olhos abertos
        \item Triste: :( com lágrimas
        \item Surpreso: o\_o com boca aberta
        \item Pensativo: :| com olho fechado (piscada)
        \item Apaixonado: <3 com corações ao redor
    \end{itemize}
    \item \textbf{Ciclo de Mudança}: cada toque avança para a próxima expressão (ciclagem)
    \item \textbf{Debounce}: software debouncing de 300ms para evitar múltiplos triggers
    \item \textbf{Benefício Distribuído}: interface local totalmente independente da API, demonstração de interação homem-máquina em sistema distribuído
    \item \textbf{Futuro}: enviar estado de interação para servidor para logging e análise de comportamento do usuário
\end{itemize}

\chapter{Arquitetura de Dados e Comunicação}

\section{Fluxo Completo de Telemetria}

\begin{figure}[H]
\centering
\begin{tikzpicture}[
    node distance=2cm and 2.5cm,
    esp/.style={rectangle, rounded corners=3mm, draw=secondary, very thick, minimum width=3cm, minimum height=0.9cm, align=center, fill=lightgreen!70},
    server/.style={rectangle, rounded corners=3mm, draw=primary, very thick, minimum width=3cm, minimum height=0.9cm, align=center, fill=lightblue!70},
    arrow/.style={-{Latex[length=2.5mm]}, very thick, draw=graytext}
]

\node[esp] (esp1) {ESP32 (Sala)};
\node[esp, below=2.5cm of esp1] (esp2) {ESP32 (Quarto)};
\node[esp, right=3cm of esp1] (esp3) {ESP8266 (Varanda)};

\node[server, below=3.5cm of esp2] (nginx) {Nginx};
\node[server, right=2.5cm of nginx] (api) {API Node.js};
\node[server, right=2.5cm of api] (db) {PostgreSQL};

\draw[arrow] (esp1) -- (nginx) node[midway, above left, font=\small]{WiFi POST};
\draw[arrow] (esp2) -- (nginx) node[midway, left, font=\small]{WiFi POST};
\draw[arrow] (esp3) -- (nginx) node[midway, above right, font=\small]{WiFi POST};
\draw[arrow] (nginx) -- (api) node[midway, above, font=\small]{HTTP};
\draw[arrow] (api) -- (db) node[midway, above, font=\small]{TCP};

\end{tikzpicture}
\caption{Fluxo de comunicação do sistema IoT distribuído.}
\end{figure}

\section{Payload JSON - Exemplo Real}

\lstset{language=Pascal}
\begin{lstlisting}
{
  "device_id": "ESP32_SALA_001",
  "device_name": "Sensor Sala Principal",
  "location": "Sala",
  "timestamp": 1715179200,
  "sensor_data": {
    "temperature_c": 24.5,
    "humidity_percent": 65.2,
    "pressure_hpa": 1013.5,
    "dht11_temperature": 24.3,
    "dht11_humidity": 64.8,
    "rssi_dbm": -45,
    "battery_voltage": 3.3,
    "uptime_seconds": 864000
  },
  "display_state": {
    "current_expression": "happy",
    "screen_active": true,
    "brightness": 255
  },
  "system_health": {
    "wifi_reconnects": 2,
    "api_failures": 0,
    "buffer_size": 3
  }
}
\end{lstlisting}

\section{Banco de Dados PostgreSQL - Schema}

A tabela no PostgreSQL (Raspberry Pi) armazena:

\begin{itemize}[leftmargin=1.2cm]
    \item \textbf{telemetry} table (séries temporais com TimescaleDB):
    \begin{itemize}
        \item \texttt{time} (TIMESTAMP com timezone)
        \item \texttt{device\_id} (VARCHAR, chave de particionamento)
        \item \texttt{temperature\_c}, \texttt{humidity\_percent}, \texttt{pressure\_hpa} (FLOAT)
        \item \texttt{rssi\_dbm} (INTEGER)
        \item \texttt{battery\_voltage} (FLOAT)
        \item índices por device\_id e time para query rápida
        \item retenção automática de 90 dias (configurável)
    \end{itemize}
    \item \textbf{devices} table:
    \begin{itemize}
        \item \texttt{id} (UUID PRIMARY KEY)
        \item \texttt{device\_id}, \texttt{name}, \texttt{location}, \texttt{status}
        \item \texttt{last\_seen} (TIMESTAMP)
        \item \texttt{firmware\_version}
        \item \texttt{ip\_address} (último IP visto)
    \end{itemize}
    \item \textbf{alerts} table:
    \begin{itemize}
        \item \texttt{id} (UUID PRIMARY KEY)
        \item \texttt{device\_id}, \texttt{alert\_type}, \texttt{severity}
        \item \texttt{created\_at}, \texttt{resolved\_at}
    \end{itemize}
\end{itemize}

\chapter{Metas de Sistemas Distribuídos}

\section{Escalabilidade}
\textbf{Objetivo}: Adicionar mais nós (novos ESP) sem redesenhar a infraestrutura. Nginx facilita escalabilidade horizontal.

\textbf{Implementação Atual}: Novos ESP podem ser adicionados e configurados com WiFi automaticamente. Nginx permanece como ponto único de entrada. Futuro: replicação de BD ou sharding por localização.

\section{Disponibilidade}
\textbf{Objetivo}: Manter o sistema funcional mesmo com falha de um ou mais componentes.

\textbf{Implementação Atual}: Buffer local em ESP permite coleta offline. Nginx e API em servidor redundante. PostgreSQL com snapshots automáticos. Futuro: replicação de BD, failover automático, standby Raspberry Pi.

\section{Tolerância a Falhas}
\textbf{Objetivo}: Reduzir impacto de desconexões WiFi, falhas de sensor, queda de serviço.

\textbf{Implementação Atual}: Reconexão automática em ESP. Retry com backoff exponencial. Health check periodicamente. Fila local de mensagens não entregues. Redundância de sensores (BME280 + DHT11).

\section{Compartilhamento de Recursos}
\textbf{Objetivo}: Coordenar acesso a sensores e dados compartilhados.

\textbf{Implementação Atual}: Nginx como ponto centralizado de roteamento. PostgreSQL como estado compartilhado único. APIs com transações ACID. Locking no banco de dados para modificações críticas.

\section{Segurança Básica}
\textbf{Objetivo}: Evitar acesso indevido a dados e controle do sistema.

\textbf{Implementação Atual}: Autenticação simples (token em sessão na web). Validação de entrada no Nginx e API. Firewalling de porta 5432 apenas para Notebook. TLS e autenticação mútua como futuro.

\chapter{Topologia Física e Geográfica}

\section{Localização dos Nós na Casa}

\begin{figure}[H]
\centering
\begin{tikzpicture}[scale=0.85]
  % Rooms
  \draw[thick, black] (0,0) rectangle (5,4);
  \draw[thick, black] (5,0) rectangle (10,4);
  \draw[thick, black] (0,4) rectangle (5,8);
  
  % Room labels
  \node[font=\large\bfseries, black] at (2.5, 2) {Sala};
  \node[font=\large\bfseries, black] at (7.5, 2) {Cozinha};
  \node[font=\large\bfseries, black] at (2.5, 6) {Quarto};
  
  % ESP nodes
  \filldraw[fill=green!40, draw=green!80, line width=2] (1.5, 1.5) circle (0.5) node[font=\small\bfseries, black] {ESP1};
  \filldraw[fill=green!40, draw=green!80, line width=2] (8, 3) circle (0.5) node[font=\small\bfseries, black] {ESP3};
  \filldraw[fill=green!40, draw=green!80, line width=2] (1, 6.5) circle (0.5) node[font=\small\bfseries, black] {ESP2};
  
  % Central servers (notebook)
  \filldraw[fill=blue!30, draw=blue!80, line width=2] (5, 6) circle (0.6) node[font=\small\bfseries, black] {NB};
  
  % Raspberry
  \filldraw[fill=blue!30, draw=blue!80, line width=2] (9, 6) circle (0.6) node[font=\small\bfseries, black] {RPI};
  
  % Router
  \filldraw[fill=orange!50, draw=orange!90, line width=2] (5, 0.5) circle (0.4) node[font=\tiny\bfseries, black] {RTR};
  
  % Antenna indicators
  \draw[dashed, gray] (5, 0.5) circle (3.5);
\end{tikzpicture}
\caption{Layout físico da infraestrutura doméstica. ESP1-3: sensores com WiFi; NB: Notebook; RPI: Raspberry Pi; RTR: Roteador WiFi (linha pontilhada mostra aproximado alcance WiFi).}
\end{figure}

\chapter{Próximas Etapas}

\section{Roadmap de Desenvolvimento}

\begin{enumerate}[leftmargin=1.2cm]
    \item \textbf{Etapa 2 (Concorrência)}: Threads em Node.js, processamento paralelo de telemetria
    \item \textbf{Etapa 3 (Sockets)}: UDP para comunicação de baixa latência entre ESP's
    \item \textbf{Etapa 4 (RPC)}: Chamadas remotas entre servidores para operações críticas
    \item \textbf{Etapa 5 (Message Bus)}: RabbitMQ ou MQTT entre servidores para desacoplamento
    \item \textbf{Etapa 6 (Service Discovery)}: mDNS ou Consul para localização dinâmica de serviços
    \item \textbf{Etapa 7 (Sincronização)}: Mutual exclusion, eleição de líder para operações críticas
    \item \textbf{Etapa 8 (Replicação)}: Espelhos de BD com consistência eventual
    \item \textbf{Etapa 9 (Resiliência)}: Circuit breaker, retry policies, graceful degradation
    \item \textbf{Etapa 10 (Segurança)}: TLS, autenticação mútua, autorização com JWT
\end{enumerate}

\chapter{Conclusão}

Este projeto demonstra um sistema distribuído genuinamente prático, rodando em infraestrutura real doméstica, com múltiplos nós heterogêneos, rede WiFi com falhas esperadas, e persistência centralizada. A infraestrutura foi cuidadosamente projetada para evitar problemas comuns (IP dinâmico, perda de WiFi) através de fixação de IPs no roteador e buffer local nos nós. Cada etapa subsequente incorporará conceitos clássicos de sistemas distribuídos de forma natural, construindo sobre essa base sólida e realista.

\end{document}
